% =============================================================================
%  Chapter 2 -- Background and Related Work
% =============================================================================
\chapter{Background and Related Work}
\label{sec:background}

\section{Citing}
\label{sec:citing}

% CITATIONS. This template uses biblatex with the authoryear style. There are
% exactly two commands you need, and the difference is grammatical:
%
%   \textcite{key}   -> "Vaswani et al. (2017) showed ..."
%                       The authors are the subject of your sentence.
%   \parencite{key}  -> "... as has been shown (Vaswani et al. 2017)."
%                       The citation is an aside.
%
% Do NOT use plain \cite: under authoryear its output is ambiguous and it will
% quietly produce a different shape than the two above.
%
% Page-specific references take an optional argument:
%   \parencite[p.~42]{key}     \textcite[see][ch.~3]{key}
%
% Add entries to bibliography/references.bib. Prefer the publisher's or DBLP's
% BibTeX over Google Scholar's, which is frequently wrong about venue and year.

Whenever a \emph{claim} is made, it must either be proven in the manuscript or
attributed to published work that supports it. An unsupported claim is the most
common thing an examiner marks in a draft.

A claim proven in the manuscript still needs a pointer to its proof, and the
reader should never have to go looking. Where the evidence sits in the same
section, say so in place, e.g., \emph{the derivation below establishes this}.
Where it sits anywhere else, cross-reference it: \Cref{sec:experiments}, never
\emph{as shown earlier}. Prefer an explicit \texttt{\textbackslash Cref} to a
bare \emph{above} or \emph{below} even within a single section, because floats
move: a figure that migrates to the next page turns \emph{the figure above}
into a small lie, and nothing warns you when it happens.

The attribution takes one of two forms, and the difference is grammatical.
% Note the \texttt{\textbackslash ...} spelling: writing a bare \textcite in
% prose makes LaTeX swallow the next character as its citation key.
Use \texttt{\textbackslash textcite} when the authors are the subject of your
sentence: as \textcite{sifa2019towards} show, financial auditors can use
machine learning tools to lighten their workload. Use
\texttt{\textbackslash parencite} when the citation is an aside, supporting a
claim you make in your own voice: machine learning is already applied to audit
workflows \parencite{sifa2019towards}.

Cite a specific page or chapter with the optional argument,
\parencite[p.~42]{example2020}, and give several sources at once by separating
the keys with commas \parencite{example2020,example2021}.

Keep the entries themselves short. An entry needs only what identifies the work, i.e.,
author, title, venue, year , and the venue should be abbreviated:
\emph{Proc.\ NeurIPS}, not \emph{Proceedings of the 30th International
Conference on Neural Information Processing Systems}. Drop pages, publisher,
address, editors, month and DOI unless a reader would genuinely struggle to
find the work without them. Exporters are the problem here: Google Scholar and
publisher sites emit every field they hold, and pasting that in unedited gives
a bibliography three times longer than it needs to be, in which the useful
information is harder to see. Trim on the way in. The
\texttt{vaswani2017attention} \parencite{vaswani2017attention} entry in \texttt{bibliography/references.bib} is
the model: four fields, and nothing a reader does not need.

\section{Figures}
\label{sec:figures}

% FIGURES. Diagram sources live in figures/ as BARE FRAGMENTS -- a file that
% begins with \begin{tikzpicture} and has no \documentclass and no
% \begin{document}. The figure environment, caption and label are supplied here
% at the call site, never inside the fragment.
%
% Why: the same fragment can then be reused in a paper or slide deck unchanged,
% with a different caption.
%
% figures/ has two subdirectories:
%   figures/template/  the logos, which stay
%   figures/examples/  the worked example diagrams, which are demonstrations
%                      and should be deleted as your own figures replace them.
%
% Placement is covered in the prose below -- it is a rule students break
% constantly, so it belongs in the rendered document, not in a comment.
%
% Captions: a full sentence ending in a period, below the figure. It should be
% understandable without the surrounding text. Every figure MUST be referenced
% from the prose -- an unreferenced figure is a layout accident.
%
% Colour: the class defines Paul Tol's `vibrant' scheme as tolorange, tolblue,
% tolcyan, tolmagenta, tolred, tolteal and tolgrey. Take them from the top of
% that list for any colour in a figure, decorative or not. The document accents
% are drawn from the same scheme (accent = tolblue, highlight = tolorange), so
% using either in a figure is safe -- but name the tol* colour rather than the
% accent, so the figure keeps its palette if the accents are ever re-pointed.

Figures are \emph{floats}: LaTeX decides where to put them. Written as a plain
\texttt{\textbackslash begin\{figure\}}, with no optional argument, a figure is
placed at the top or bottom of a page, or on a page of its own, wherever it
fits without tearing a hole in the running text. That default is almost always
the right answer, and it is the reason figures are numbered: the reader is sent
to them by number, not by position.

Placement can be forced with an optional argument, and this is where most
figure trouble starts. \texttt{[h]} asks for \emph{here}; \texttt{[!h]} asks
more insistently, telling LaTeX to ignore its own limits on how much of a page
a float may occupy; \texttt{[H]}, from the \texttt{float} package, nails the
figure in place and refuses to move it at all. Each degrades the page when the
figure does not fit: \texttt{[!h]} commonly defers it to the next page anyway,
having first left a short page behind, and \texttt{[H]} opens a large area of
white space above the figure rather than reflowing anything.

Reserve forced placement for the cases that genuinely need it, such as a run of
small diagrams that must be read in the order they appear, or a figure that a
sentence physically points at. Everywhere else, leave the option off.
\Cref{fig:neural_network} carries no placement option, and LaTeX has put it
wherever it fitted; the sentence you are reading does not care where that was.
It follows that the prose should never say \emph{the figure below}: name the
figure, and the sentence stays true wherever the figure lands.

Give every caption a short form as well, as
\texttt{\textbackslash caption[short]\{long\}}: the short one is what the List
of Figures prints, which keeps that list scannable instead of reproducing three
lines of explanation per entry.

A figure taken or adapted from elsewhere must say so in its caption, with a
citation like any other borrowed material; \Cref{fig:neural_network} shows an example how to do this.

For neural network architectures in particular, we recommend the
NNTikZ\footnote{\href{https://github.com/fraserlove/nntikz}{github.com/fraserlove/nntikz}.}
collection \parencite{love2024nntikz} as a starting point.


\begin{figure}
\centering
% =============================================================================
%  Fully-connected neural network.
%
%  Adapted from NNTikZ, Copyright (c) 2024 Fraser Love, MIT Licence:
%  https://github.com/fraserlove/nntikz -- re-coloured with this template's
%  palette and reduced to a bare fragment.
%
%  Colours come from the Tol `vibrant' scheme defined in aml-thesis.cls:
%  input tolblue, hidden tolmagenta, output tolorange. Wiring and layer boxes
%  are black -- they are structure, not data, so they take no palette colour.
%  Nothing here hard-codes a hex value, so the figure follows the palette if it
%  is ever re-pointed.
% =============================================================================
\begin{tikzpicture}[
    >=LaTeX,
    node/.style={circle, minimum width=2.25em, draw, thick},
    innode/.style={node, draw=tolblue,   fill=tolblue!12},
    hidnode/.style={node, draw=tolmagenta, fill=tolmagenta!12},
    outnode/.style={node, draw=tolorange, fill=tolorange!12},
    cell/.style={rectangle, rounded corners=2mm, minimum height=2.5em,
                 minimum width=2.5em, draw=black, thick, dashed},
    arrow/.style={-latex, semithick, draw=black}
]

    % Input layer
    \foreach \x in {1,...,2}
        \draw node at (0, -\x*1.25 - 0.625) [innode] (first_\x) {$x_\x$};
    \draw node at (0, -5*1.25 + 0.625) [innode] (first_n) {$x_n$};
    \path (first_2) -- (first_n) node[pos=0.38, scale=2] {\vdots};

    % First hidden layer
    \foreach \x in {1,...,3}
        \node at (2.5, -\x*1.25) [hidnode] (second_\x){};
    \draw node at (2.5, -5*1.25) [hidnode] (second_m) {};
    \path (second_3) -- (second_m) node[pos=0.38, scale=2] {\vdots};

    % Second hidden layer
    \foreach \x in {1,...,3}
        \node at (5, -\x*1.25) [hidnode] (third_\x){};
    \draw node at (5, -5*1.25) [hidnode] (third_m) {};
    \path (third_3) -- (third_m) node[pos=0.38, scale=2] {\vdots};

    % Output layer
    \foreach \x in {1,...,2}
        \node at (7.5, -\x*1.25 - 0.625) [outnode] (fourth_\x){$\hat{y}_\x$};
    \draw node at (7.5, -5*1.25 + 0.625) [outnode] (fourth_k) {$\hat{y}_m$};
    \path (fourth_2) -- (fourth_k) node[pos=0.38, scale=2] {\vdots};

    % Input -> hidden 1
    \foreach \i in {1,...,2}
        \foreach \j in {1,...,3}
            \draw [arrow] (first_\i) to (second_\j);
    \foreach \i in {1,...,2}
        \draw [arrow] (first_\i) to (second_m);
    \foreach \i in {1,...,3}
        \draw [arrow] (first_n) to (second_\i);
    \draw [arrow] (first_n) to (second_m);

    % Hidden 1 -> hidden 2
    \foreach \i in {1,...,3}
        \foreach \j in {1,...,3}
            \draw [arrow] (second_\i) to (third_\j);
    \foreach \i in {1,...,3}
        \draw [arrow] (second_\i) to (third_m);
    \foreach \i in {1,...,3}
        \draw [arrow] (second_m) to (third_\i);
    \draw [arrow] (second_m) to (third_m);

    % Hidden 2 -> output
    \foreach \i in {1,...,3}
        \foreach \j in {1,...,2}
            \draw [arrow] (third_\i) to (fourth_\j);
    \foreach \i in {1,...,3}
        \draw [arrow] (third_\i) to (fourth_k);
    \foreach \i in {1,...,2}
        \draw [arrow] (third_m) to (fourth_\i);
    \draw [arrow] (third_m) to (fourth_k);

    % Layer labels
    \node[above=0.25cm of first_1]  {$\bm{x}$};
    \node[above=0.25cm of second_1] {$\bm{h}^{(1)}$};
    \node[above=0.25cm of third_1]  {$\bm{h}^{(2)}$};
    \node[above=0.25cm of fourth_1] {$\hat{\bm{y}}$};

    % Layer groupings
    \node[cell, fit=(first_1) (first_n)]   {};
    \node[cell, fit=(second_1) (second_m)] {};
    \node[cell, fit=(third_1) (third_m)]   {};
    \node[cell, fit=(fourth_1) (fourth_k)] {};
\end{tikzpicture}

\caption[A fully-connected feed-forward network]{A fully-connected
feed-forward network mapping an input $\bm{x}$ to a
prediction $\hat{\bm{y}}$ through two hidden layers. A caption is a full
sentence and explains what the reader is looking at. Adapted from
\textcite{love2024nntikz}.}
\label{fig:neural_network}
\end{figure}

\subsection{Colour}
\label{sec:colour}

Whenever colour is used at all, it should come from a colour-blind-safe
palette.  Red--green colour vision deficiency affects roughly
one in twelve men and one in 250 women of European descent, and between
4\% and 6.5\% of men of East Asian descent \parencite{birch2012worldwide}.
A figure that distinguishes its series by red against green is therefore
unreadable to part of any audience, and a thesis is read by a small audience
in which one such reader is enough.

Which palette is a matter of taste: any scheme designed for colour vision
deficiency will serve, and several good ones exist. As a starting point we
suggest Paul Tol's schemes,\footnote{Paul Tol, \emph{Colour Schemes}:
\href{https://sronpersonalpages.nl/~pault/}{sronpersonalpages.nl/\textasciitilde pault}.}
of which this template ships the \emph{vibrant} one as a worked example,
defined in \texttt{aml-thesis.cls} as
\textcolor{tolorange}{\texttt{tolorange}},
\textcolor{tolblue}{\texttt{tolblue}},
\textcolor{tolcyan}{\texttt{tolcyan}},
\textcolor{tolmagenta}{\texttt{tolmagenta}},
\textcolor{tolred}{\texttt{tolred}},
\textcolor{tolteal}{\texttt{tolteal}} and
\textcolor{tolgrey}{\texttt{tolgrey}}. \Cref{fig:neural_network,fig:network_pair}
are coloured this way, showing the palette in action. Both are deliberately
more colourful than they likely need to be, in order to display the scheme; colour
figures to your own taste.

For figures generated in Python, Tol's schemes are packaged as
\texttt{tol\_colors},\footnote{\href{https://github.com/Descanonge/tol_colors}{github.com/Descanonge/tol\_colors}.}
which plugs into Matplotlib directly, so plots match the diagrams drawn in
TikZ.

\subsection{Sub-figures}
\label{sec:subfigures}

% Two related panels side by side. Each subfigure gets its own caption and
% label; reference them as \cref{fig:neuron} or \cref{fig:network_pair}, and
% \subref gives just the (a)/(b) marker for use inside the main caption.
% The width arguments should sum to slightly less than 1.0 to leave a gutter.

Two panels that belong together should share one figure number rather than
take two. \Cref{fig:network_pair} does this: each panel carries its own
caption and label, so \cref{fig:neuron} and \cref{fig:rnn} can be cited
separately, while the shared caption explains what the pair has in common.

\begin{figure}
\centering
% Each panel is boxed to a common height and centred inside it, so the two
% sit on a shared vertical axis and their captions line up. Without the
% minipage the shorter panel hangs at the bottom of the taller one.
\begin{subfigure}[b]{0.44\textwidth}
  \centering
  \begin{minipage}[c][45mm][c]{\linewidth}
    \centering
    \adjustbox{max width=\linewidth, max totalheight=45mm}{%
      % =============================================================================
%  A single artificial neuron.
%
%  Adapted from NNTikZ, Copyright (c) 2024 Fraser Love, MIT Licence:
%  https://github.com/fraserlove/nntikz -- re-coloured with this template's
%  palette and reduced to a bare fragment.
%
%  Roles match the other network figures: inputs tolblue, unit tolmagenta,
%  output tolorange, wiring black.
% =============================================================================
\begin{tikzpicture}[
    >=LaTeX,
    node/.style={circle, minimum width=2.25em, draw, thick},
    innode/.style={node, draw=tolblue,     fill=tolblue!12},
    midnode/.style={node, draw=tolmagenta, fill=tolmagenta!12},
    outnode/.style={node, draw=tolorange,  fill=tolorange!12},
    arrow/.style={-latex, semithick, draw=black},
    wlabel/.style={midway, fill=white, inner sep=1.5pt}
]

    % Inputs
    \foreach \x in {0,1,2}
        \node[innode] (x\x) at (0,-\x*1.25) {$x_{\x}$};
    \node[innode] (xm) at (0,-4*1.25) {$x_{n}$};
    \path (x2) -- (xm) node[pos=0.35, scale=2] {\vdots};

    % Summation unit and output
    \node[midnode] (z) at (2.5,-2.5) {$z$};
    \node[outnode, right of=z, node distance=2.5cm] (yhat) {$\hat{y}$};

    % Weighted inputs
    \draw[arrow] (x0) -- (z) node[wlabel] {$w_0$};
    \foreach \x in {1,...,2}
        \draw[arrow] (x\x) -- (z) node[wlabel] {$w_{\x}$};
    \draw[arrow] (xm) -- (z) node[wlabel] {$w_{n}$};

    % Activation
    \draw[arrow] (z) -- (yhat) node[wlabel] {$\sigma(z)$};

    \node[above of=z, node distance=1.2cm] {$z = \bm{w}^\top \bm{x}$};
\end{tikzpicture}
}
  \end{minipage}
  \caption{A single unit.}
  \label{fig:neuron}
\end{subfigure}
\hfill
\begin{subfigure}[b]{0.52\textwidth}
  \centering
  \begin{minipage}[c][45mm][c]{\linewidth}
    \centering
    \adjustbox{max width=\linewidth, max totalheight=45mm}{%
      % =============================================================================
%  A recurrent cell, folded and unrolled over time.
%
%  Adapted from NNTikZ, Copyright (c) 2024 Fraser Love, MIT Licence:
%  https://github.com/fraserlove/nntikz -- re-coloured with this template's
%  palette and reduced to a bare fragment.
%
%  Roles match the other network figures: inputs tolblue, recurrent cells
%  tolmagenta, outputs tolorange, wiring black. Needs the `chains' TikZ
%  library, loaded by aml-thesis.cls.
% =============================================================================
\begin{tikzpicture}[
    >=LaTeX,
    cell/.style={rectangle, rounded corners=2mm, minimum height=2.5em,
                 minimum width=2.5em, draw=tolmagenta, fill=tolmagenta!12, thick},
    cellc/.style={cell, on chain, join},
    innode/.style={circle, minimum width=2.25em, draw=tolblue,
                   fill=tolblue!12, thick},
    outnode/.style={circle, minimum width=2.25em, draw=tolorange,
                    fill=tolorange!12, thick},
    arrow/.style={-latex, semithick, draw=black}
]

    % Unrolled over time
    \begin{scope}[start chain=going right, nodes=cellc, arrow,
                  local bounding box=chain]
        \path[shift={(-9em, 1em)}]
            node (h1) {$\bm{h}_1$}
            node     (h2) {$\bm{h}_2$}
            node[xshift=2em] (hT) {$\bm{h}_{T}$};
    \end{scope}

    \foreach \X in {1, 2, T} {
        \draw[arrow, latex-] (h\X.south) -- ++ (0, -2em)
            node[below, innode] (x\X) {$\bm{x}_\X$};
    }
    \foreach \X in {1, 2, T} {
        \draw[arrow] (h\X.north) -- ++ (0, 2em)
            node[above, outnode] (y\X) {$\hat{\bm{y}}_\X$};
    }
    \path (x2) -- (xT) node[midway, scale=2] {\dots};

    \node[left=1em of chain, scale=2] (eq) {$=$};

    % Folded form
    \node[cell, left=2em of eq] (AL) {$\bm{h}_t$};
    \path (AL.west) ++ (-1em,2em) coordinate (aux);
    \draw[arrow, rounded corners] (AL.east) -| ++ (1em, 2em) -- (aux) |- (AL.west);
    % White mask so the recurrent loop passes behind the outgoing arrow.
    \draw[white, line width=0.5em] (AL.north) -- ++ (0, 1.9em);
    \draw[arrow] (AL.north) -- ++ (0, 2em) node[above, outnode] {$\hat{\bm{y}}_t$};
    \draw[arrow, latex-] (AL.south) -- ++ (0, -2em) node[below, innode] {$\bm{x}_t$};
\end{tikzpicture}
}
  \end{minipage}
  \caption{A recurrent cell, folded and unrolled.}
  \label{fig:rnn}
\end{subfigure}
\caption[A single unit and a recurrent cell]{Two views of the same machinery:
(\subref{fig:neuron})~one unit,
summing weighted inputs and applying an activation; (\subref{fig:rnn})~a
recurrent cell shown folded and unrolled over $T$ steps. Each sub-figure takes
its own caption and label, and \texttt{adjustbox} scales an oversized fragment
rather than its coordinates being edited. Adapted from
\textcite{love2024nntikz}.}
\label{fig:network_pair}
\end{figure}
